% Ce fichier est prévu pour être inclus dans un document LaTeX à deux colonnes.
% Packages requis dans le préambule :
% \usepackage{amsmath,amssymb,array,booktabs,tabularx}

\section{Évaluation de GEAR avec des participants}
\label{sec:evaluation-participants}

L'évaluation technique présentée dans le tool track vérifie que GEAR
couvre les fonctionnalités des frameworks étudiés et génère des
implémentations exécutables qui préservent la sémantique des workflows.
Cette évaluation ne permet cependant pas de déterminer si GEAR aide
effectivement les développeurs à concevoir, modifier et porter des
systèmes multi-agents.

Nous conduisons donc une étude contrôlée avec des participants afin de
comparer deux approches de développement :

\begin{itemize}
    \item \textbf{Condition GEAR :} le participant spécifie le système
    dans GEAR, valide sa configuration, puis génère les implémentations
    correspondant aux frameworks cibles ;

    \item \textbf{Condition native :} le participant développe
    directement le système en utilisant les abstractions, l'API et les
    mécanismes d'orchestration du framework cible.
\end{itemize}

L'étude examine quatre dimensions complémentaires : l'efficacité,
l'efficience, la maintenance et la portabilité, ainsi que l'expérience
utilisateur.


\subsection{Questions de recherche}
\label{subsec:questions-recherche-participants}

L'étude est guidée par les questions de recherche suivantes.

\begin{quote}
\textbf{RQ1 :} Dans quelle mesure GEAR aide-t-il les développeurs à
produire des systèmes multi-agents corrects ?
\end{quote}

Cette question porte sur l'efficacité de GEAR. Elle étudie la capacité
des participants à produire des systèmes qui respectent les exigences
et réussissent les tests fonctionnels.

\begin{quote}
\textbf{RQ2 :} Dans quelle mesure GEAR réduit-il l'effort nécessaire
pour développer un système multi-agent ?
\end{quote}

Cette question porte sur l'efficience du processus de développement.
Elle prend notamment en compte le temps, le nombre de modifications et
les tentatives nécessaires pour obtenir une solution correcte.

\begin{quote}
\textbf{RQ3 :} Dans quelle mesure GEAR facilite-t-il la modification et
la migration des systèmes multi-agents entre différents frameworks ?
\end{quote}

Cette question évalue le principal objectif de portabilité de GEAR. Elle
mesure l'effort nécessaire pour faire évoluer un système existant et
pour le migrer vers un autre framework.

\begin{quote}
\textbf{RQ4 :} Comment les développeurs perçoivent-ils l'utilisabilité,
la charge de travail et l'utilité de GEAR ?
\end{quote}

Cette question étudie l'expérience des participants, leur confiance
dans les artefacts générés, ainsi que les avantages et limites perçus
de GEAR.


\subsection{Hypothèses}
\label{subsec:hypotheses-participants}

À partir des objectifs de GEAR, nous formulons les hypothèses suivantes :

\begin{description}
    \item[\textbf{H1 :}]
    les participants utilisant GEAR produisent davantage de systèmes
    corrects que ceux utilisant directement les frameworks natifs ;

    \item[\textbf{H2 :}]
    les participants utilisant GEAR terminent les tâches de
    développement plus rapidement ;

    \item[\textbf{H3 :}]
    GEAR réduit le nombre d'erreurs et de tentatives nécessaires pour
    obtenir une solution correcte ;

    \item[\textbf{H4 :}]
    GEAR réduit l'effort nécessaire pour modifier un système
    multi-agent existant ;

    \item[\textbf{H5 :}]
    GEAR réduit l'effort nécessaire pour migrer un système multi-agent
    vers un autre framework ;

    \item[\textbf{H6 :}]
    les participants utilisant GEAR rapportent une utilisabilité plus
    élevée et une charge de travail plus faible.
\end{description}


\subsection{Plan expérimental}
\label{subsec:plan-experimental}

Nous adoptons un plan expérimental intra-sujets. Chaque participant
réalise une série de tâches avec GEAR et une autre série de tâches avec
les frameworks natifs. Cette approche permet de comparer les deux
conditions chez un même participant et de limiter l'influence des
différences individuelles.

L'ordre des conditions est contrebalancé afin de limiter les effets
d'apprentissage et de fatigue. Deux ensembles de tâches, notés $A$ et
$B$, sont utilisés. Ils portent sur des domaines différents, mais
présentent une complexité équivalente en matière de nombre d'agents,
de dépendances et de structures d'orchestration.

Une répartition simplifiée des participants est présentée dans le
Tableau~\ref{tab:contrebalancement}.

\begin{table}[t]
    \centering
    \caption{Contrebalancement des conditions expérimentales.}
    \label{tab:contrebalancement}
    \small
    \begin{tabularx}{\columnwidth}{@{}l>{\raggedright\arraybackslash}X>{\raggedright\arraybackslash}X@{}}
        \toprule
        \textbf{Groupe}
        & \textbf{Première condition}
        & \textbf{Deuxième condition} \\
        \midrule
        Groupe 1
            & GEAR, tâche A
            & Native, tâche B \\
        Groupe 2
            & Native, tâche A
            & GEAR, tâche B \\
        \bottomrule
    \end{tabularx}
\end{table}

L'affectation des participants aux groupes est réalisée aléatoirement
tout en maintenant des groupes de tailles comparables. La direction de
migration entre les frameworks est également équilibrée entre les
participants.


\subsection{Frameworks utilisés}
\label{subsec:frameworks-experience}

Il n'est pas réaliste de demander aux participants de manipuler les dix
frameworks pris en charge par GEAR. Nous sélectionnons donc deux
frameworks représentant des paradigmes suffisamment différents :

\begin{itemize}
    \item \textbf{CrewAI}, qui organise principalement les systèmes
    autour des rôles, des agents et des tâches ;

    \item \textbf{LangGraph}, qui représente les workflows sous la
    forme de graphes avec gestion d'un état.
\end{itemize}

Ces deux frameworks permettent d'évaluer la capacité de GEAR à masquer
les différences entre une approche orientée rôles et tâches et une
approche orientée graphes.

Le choix de deux frameworks ne vise pas à réévaluer la couverture
technique de GEAR, celle-ci étant étudiée séparément sur les dix
connecteurs. Il vise à mesurer le bénéfice humain apporté par
l'abstraction indépendante des frameworks.


\subsection{Tâches expérimentales}
\label{subsec:taches-experimentales}

Chaque condition comporte trois types de tâches : construction,
modification et migration.


\subsubsection{Tâche de construction}

Le participant doit construire un système multi-agent séquentiel
comportant plusieurs agents :

\begin{equation}
A \rightarrow B \rightarrow C.
\end{equation}

Cette tâche nécessite notamment de :

\begin{itemize}
    \item définir les agents et leurs responsabilités ;
    \item définir les tâches associées ;
    \item configurer les modèles et les paramètres nécessaires ;
    \item définir les dépendances ;
    \item produire une implémentation exécutable.
\end{itemize}

Cette tâche mesure la facilité de prise en main, la compréhension des
concepts et la capacité à construire un premier système fonctionnel.


\subsubsection{Tâche de modification}

Le participant reçoit ensuite une nouvelle exigence. Le système initial
doit évoluer vers une architecture comprenant une exécution parallèle
et une synchronisation :

\begin{equation}
A \rightarrow (B \parallel C) \rightarrow D.
\end{equation}

Le participant doit :

\begin{itemize}
    \item ajouter un nouvel agent ;
    \item introduire une exécution parallèle ;
    \item configurer une synchronisation fan-in ;
    \item préserver la transmission des résultats ;
    \item maintenir le comportement déjà validé.
\end{itemize}

Cette tâche mesure l'effort de maintenance et la capacité à faire
évoluer un système existant.


\subsubsection{Tâche de migration}

Enfin, le participant doit migrer le système obtenu vers un second
framework.

Dans la condition native, il doit adapter ou réécrire l'implémentation
en utilisant les abstractions du nouveau framework. Dans la condition
GEAR, il conserve la spécification indépendante du framework et
sélectionne une nouvelle cible de génération.

Cette tâche mesure directement la portabilité offerte par GEAR. Elle
constitue la tâche principale pour répondre à RQ3.


\subsection{Équivalence des tâches}
\label{subsec:equivalence-taches}

Les ensembles de tâches $A$ et $B$ utilisent des scénarios fonctionnels
différents afin d'éviter que les participants reproduisent directement
une solution déjà construite. Ils conservent cependant :

\begin{itemize}
    \item le même nombre d'agents ;
    \item le même nombre de tâches ;
    \item le même nombre de dépendances ;
    \item les mêmes structures séquentielles et parallèles ;
    \item le même nombre de modifications ;
    \item des critères de réussite comparables.
\end{itemize}

L'équivalence des tâches est vérifiée au cours d'une étude pilote.


\subsection{Participants}
\label{subsec:participants}

L'étude cible des participants possédant des connaissances de base en
programmation Python et une compréhension générale des \gls{llm} et des
agents logiciels. Les participants peuvent être :

\begin{itemize}
    \item des doctorants ;
    \item des chercheurs ;
    \item des ingénieurs de recherche ;
    \item des développeurs ;
    \item des praticiens travaillant sur des systèmes agentiques.
\end{itemize}

Avant l'expérience, nous collectons les informations suivantes :

\begin{itemize}
    \item expérience générale en programmation ;
    \item expérience en Python ;
    \item connaissance des \gls{llm} ;
    \item expérience avec les systèmes multi-agents ;
    \item expérience avec CrewAI ;
    \item expérience avec LangGraph ;
    \item expérience avec les outils de modélisation ;
    \item expérience avec les lignes de produits logiciels.
\end{itemize}

Ces informations sont utilisées pour décrire l'échantillon et contrôler
l'influence de l'expérience préalable sur les résultats.


\subsubsection{Nombre de participants}

Le nombre définitif de participants est déterminé au moyen d'une
analyse de puissance statistique réalisée à partir des résultats de
l'étude pilote.

À titre indicatif, pour détecter un effet moyen dans une comparaison
intra-sujets avec une puissance de $0{,}80$ et un seuil
$\alpha=0{,}05$, nous visons environ 34 participants disposant de
données exploitables. Nous pouvons ainsi recruter entre 38 et 40
participants afin d'anticiper les abandons, les problèmes techniques et
les données incomplètes.

Si ce nombre ne peut pas être atteint, les résultats devront être
présentés comme exploratoires et accompagnés d'intervalles de confiance
et de tailles d'effet.


\subsection{Déroulement de l'expérience}
\label{subsec:deroulement-experience}

Une session expérimentale dure approximativement entre 90 et
120 minutes. Son déroulement est présenté dans le
Tableau~\ref{tab:deroulement-session}.

\begin{table}[t]
    \centering
    \caption{Déroulement indicatif d'une session expérimentale.}
    \label{tab:deroulement-session}
    \begin{tabular}{@{}lr@{}}
        \toprule
        \textbf{Phase} & \textbf{Durée} \\
        \midrule
        Information et consentement
            & 5 min \\
        Questionnaire démographique
            & 5 min \\
        Tutoriel GEAR
            & 15 min \\
        Présentation de l'environnement
            & 10 min \\
        Première condition
            & 20--25 min \\
        Questionnaire intermédiaire
            & 5 min \\
        Pause
            & 5 min \\
        Deuxième condition
            & 20--25 min \\
        Questionnaire final
            & 10 min \\
        Entretien semi-directif
            & 10--15 min \\
        \bottomrule
    \end{tabular}
\end{table}

Le temps consacré au tutoriel n'est pas inclus dans le temps de
réalisation des tâches expérimentales. Chaque participant reçoit les
mêmes explications, exemples et ressources.


\subsection{Variables et mesures}
\label{subsec:variables-mesures}

La variable indépendante principale est la méthode de développement :

\begin{equation}
\text{Condition}
\in
\{\text{GEAR},\text{Native}\}.
\end{equation}

Les principales variables dépendantes sont présentées dans le
Tableau~\ref{tab:mesures-participants}.

\begin{table*}[t]
    \centering
    \caption{Mesures utilisées dans l'étude avec participants.}
    \label{tab:mesures-participants}
    \begin{tabularx}{\textwidth}{@{}l>{\raggedright\arraybackslash}X>{\raggedright\arraybackslash}X@{}}
        \toprule
        \textbf{Dimension}
        & \textbf{Mesures objectives}
        & \textbf{Mesures subjectives} \\
        \midrule
        Efficacité
            & Tâches terminées, tests réussis, erreurs fonctionnelles,
              conformité du workflow
            & Confiance dans la solution produite \\
        Efficience
            & Temps, nombre de tentatives, nombre de modifications,
              consultations de la documentation
            & Effort perçu \\
        Maintenance
            & Temps de modification, fichiers modifiés, régressions,
              erreurs introduites
            & Facilité de modification perçue \\
        Portabilité
            & Temps de migration, quantité de code modifié, erreurs de
              migration, tests réussis après migration
            & Facilité de migration perçue \\
        Expérience utilisateur
            & --- 
            & SUS, NASA-TLX, utilité perçue et intention de
              réutilisation \\
        \bottomrule
    \end{tabularx}
\end{table*}


\subsubsection{Taux de réussite}

Le taux de réussite pour la condition $c$ est défini par :

\begin{equation}
R_{\text{succès}}(c)
=
\frac{
    n_{\text{tâches terminées},c}
}{
    n_{\text{tâches assignées},c}
}.
\end{equation}


\subsubsection{Correction fonctionnelle}

Soit $n_{\mathrm{tests}}$ le nombre total de tests associés à une tâche
et $n_{\text{tests réussis}}$ le nombre de tests réussis par la
solution du participant. La correction fonctionnelle est définie par :

\begin{equation}
C_{\mathrm{fonctionnelle}}
=
\frac{
    n_{\text{tests réussis}}
}{
    n_{\mathrm{tests}}
}.
\end{equation}


\subsubsection{Temps de réalisation}

Le temps de réalisation correspond au temps écoulé entre le début de la
tâche et :

\begin{itemize}
    \item la validation complète de la solution ;
    \item l'abandon de la tâche ;
    \item l'expiration du temps maximal autorisé.
\end{itemize}

Les tâches non terminées sont conservées dans l'analyse et identifiées
comme des observations censurées ou des échecs, selon l'analyse
statistique retenue.


\subsubsection{Effort de modification}

L'effort de modification est caractérisé par :

\begin{itemize}
    \item le temps nécessaire pour satisfaire la nouvelle exigence ;
    \item le nombre de fichiers modifiés ;
    \item le nombre de lignes ajoutées, supprimées ou modifiées ;
    \item le nombre de tentatives avant la réussite ;
    \item le nombre de tests précédemment valides qui échouent après
    la modification.
\end{itemize}


\subsubsection{Effort de migration}

L'effort de migration est caractérisé par :

\begin{itemize}
    \item le temps nécessaire pour obtenir une implémentation
    fonctionnelle dans le second framework ;
    \item le nombre de fichiers et de lignes modifiés ;
    \item le nombre d'erreurs rencontrées ;
    \item le nombre de consultations de la documentation ;
    \item la proportion de tests réussis après la migration.
\end{itemize}


\subsection{Questionnaires}
\label{subsec:questionnaires}

L'utilisabilité de GEAR est évaluée avec le \emph{System Usability
Scale} (SUS)~\cite{brooke1996sus}. La charge de travail perçue après
chaque condition est évaluée avec le questionnaire
NASA-TLX~\cite{hart1988nasa}.

Afin de limiter la durée de l'expérience, ces questionnaires sont
complétés par un nombre restreint de questions portant spécifiquement
sur GEAR :

\begin{itemize}
    \item GEAR a-t-il facilité la compréhension de la structure du
    système ?
    \item Les diagnostics ont-ils aidé à identifier et corriger les
    erreurs ?
    \item La génération de code a-t-elle semblé prévisible ?
    \item Le participant ferait-il confiance au code généré ?
    \item GEAR a-t-il facilité la modification du système ?
    \item GEAR a-t-il facilité la migration vers un autre framework ?
    \item Le participant envisagerait-il d'utiliser GEAR dans un projet
    réel ?
\end{itemize}


\subsection{Entretien semi-directif}
\label{subsec:entretien-participants}

À la fin de la session, un entretien semi-directif est réalisé afin de
mieux comprendre les résultats quantitatifs. Les questions portent
notamment sur :

\begin{itemize}
    \item les difficultés rencontrées dans chaque condition ;
    \item les fonctionnalités de GEAR les plus utiles ;
    \item les fonctionnalités manquantes ;
    \item la compréhension du lien entre la spécification et le code
    généré ;
    \item la confiance dans les artefacts générés ;
    \item les contextes dans lesquels le participant utiliserait GEAR ;
    \item les limites perçues de l'approche.
\end{itemize}


\subsection{Environnement expérimental}
\label{subsec:environnement-experimental}

Tous les participants utilisent un environnement standardisé contenant :

\begin{itemize}
    \item les mêmes versions de Python et des frameworks ;
    \item les mêmes versions de GEAR ;
    \item les mêmes dépendances ;
    \item les mêmes modèles de départ ;
    \item les mêmes ressources documentaires ;
    \item les mêmes tests automatiques ;
    \item les mêmes limites de temps.
\end{itemize}

L'environnement peut être distribué sous la forme d'un conteneur afin
de réduire les différences de configuration entre les participants.

La correction des solutions est évaluée automatiquement à l'aide des
mêmes tests pour l'ensemble des participants. Les identifiants des
participants sont anonymisés avant l'analyse.


\subsection{Utilisation des outils externes}
\label{subsec:outils-externes}

L'utilisation d'outils externes doit être identique dans les deux
conditions. Deux stratégies sont possibles :

\begin{enumerate}
    \item fournir une documentation locale commune et interdire les
    assistants de programmation afin de maximiser le contrôle
    expérimental ;

    \item autoriser la documentation en ligne et les assistants de
    programmation afin de se rapprocher des pratiques réelles, tout en
    enregistrant précisément leur utilisation.
\end{enumerate}

Pour une première étude contrôlée, nous privilégions la première
stratégie. Les participants disposent d'une documentation équivalente
pour GEAR, CrewAI et LangGraph, mais ne peuvent pas utiliser
d'assistant génératif pendant les tâches.


\subsection{Contrôle des biais}
\label{subsec:controle-biais}

Plusieurs mécanismes sont utilisés pour limiter les biais.


\subsubsection{Expérience préalable}

L'expérience des participants avec Python, GEAR, CrewAI et LangGraph est
mesurée avant l'étude et intégrée comme variable de contrôle dans les
analyses.


\subsubsection{Effet d'apprentissage}

L'ordre des conditions est contrebalancé et les participants utilisent
des tâches différentes mais équivalentes dans les deux conditions.


\subsubsection{Effet de fatigue}

Une pause est proposée entre les deux conditions et l'ordre des
conditions est réparti entre les participants.


\subsubsection{Biais d'évaluation}

La correction des solutions est principalement évaluée par des tests
automatiques. Les évaluateurs ne connaissent pas l'identité des
participants lors de l'analyse des solutions.


\subsubsection{Biais lié aux tâches}

L'équivalence des tâches est vérifiée au cours du pilote à partir du
temps de réalisation, du nombre d'erreurs et des retours des
participants.


\subsection{Étude pilote}
\label{subsec:etude-pilote}

Avant l'étude principale, nous réalisons une étude pilote avec quatre à
six participants. Le pilote vise à vérifier :

\begin{itemize}
    \item la clarté des consignes ;
    \item la durée des différentes phases ;
    \item l'équivalence des ensembles de tâches ;
    \item le fonctionnement des environnements ;
    \item la fiabilité des tests automatiques ;
    \item la collecte des temps, événements et erreurs ;
    \item la compréhension de GEAR par les participants ;
    \item la pertinence des questionnaires.
\end{itemize}

Si les tâches ou le protocole sont modifiés de manière substantielle
après le pilote, les données des participants du pilote ne sont pas
intégrées à l'analyse principale.


\subsection{Analyse quantitative}
\label{subsec:analyse-quantitative}

Nous rapportons, pour chaque mesure, les statistiques descriptives, les
tailles d'effet et les intervalles de confiance à 95\,\%.

Dans la mesure où chaque participant réalise plusieurs tâches, nous
privilégions des modèles à effets mixtes. Pour une mesure continue telle
que le temps, le modèle peut être formulé de la manière suivante :

\begin{equation}
\begin{aligned}
Y \sim {} & \text{Condition} + \text{Ordre} + \text{Expérience} \\
           & + (1 \mid \text{Participant})
             + (1 \mid \text{Tâche}).
\end{aligned}
\end{equation}

La condition expérimentale est considérée comme un effet fixe, tandis
que le participant et la tâche sont considérés comme des effets
aléatoires.

Pour les résultats binaires, tels que la réussite ou l'échec d'une
tâche, nous utilisons un modèle linéaire généralisé à effets mixtes avec
une fonction de lien logistique.

Pour les variables de comptage, telles que le nombre d'erreurs, nous
utilisons un modèle de Poisson ou un modèle binomial négatif selon la
dispersion observée.

Lorsque des analyses appariées plus simples sont appropriées, nous
utilisons :

\begin{itemize}
    \item un test de Student apparié si les hypothèses paramétriques
    sont satisfaites ;
    \item un test de Wilcoxon apparié dans le cas contraire.
\end{itemize}

Les comparaisons multiples sont corrigées, par exemple avec la méthode
de Holm. Les résultats ne sont pas interprétés uniquement à partir des
valeurs $p$ : les tailles d'effet et leurs intervalles de confiance sont
également rapportés.


\subsection{Analyse qualitative}
\label{subsec:analyse-qualitative}

Les entretiens sont analysés au moyen d'une analyse thématique. Le
processus comprend :

\begin{enumerate}
    \item la transcription et l'anonymisation des entretiens ;
    \item une première lecture des données ;
    \item l'identification de codes initiaux ;
    \item le regroupement des codes en thèmes ;
    \item la comparaison des thèmes avec les résultats quantitatifs ;
    \item la sélection de verbatims représentatifs.
\end{enumerate}

Afin de renforcer la fiabilité de l'analyse, une partie des entretiens
est codée indépendamment par au moins deux chercheurs. Les divergences
sont ensuite discutées afin d'aboutir à une interprétation commune.


\subsection{Considérations éthiques et protection des données}
\label{subsec:ethique-participants}

Avant la collecte des données, le protocole est soumis à la procédure
éthique applicable dans l'établissement des chercheurs. Chaque
participant reçoit une fiche d'information précisant :

\begin{itemize}
    \item les objectifs de l'étude ;
    \item le déroulement et la durée de la session ;
    \item les données collectées ;
    \item les éventuels enregistrements d'écran ou audio ;
    \item la méthode d'anonymisation ;
    \item la durée de conservation des données ;
    \item les conditions de retrait de l'étude.
\end{itemize}

Le consentement libre et éclairé des participants est recueilli avant
le début de l'expérience. Les données sont anonymisées et traitées
conformément aux exigences applicables en matière de protection des
données, notamment le RGPD.


\subsection{Prérégistration et matériel de réplication}
\label{subsec:preregistration-replication}

Avant l'étude principale, nous préenregistrons, lorsque cela est
possible :

\begin{itemize}
    \item les questions de recherche ;
    \item les hypothèses ;
    \item le plan expérimental ;
    \item les variables étudiées ;
    \item les critères d'exclusion ;
    \item le nombre de participants visé ;
    \item les méthodes d'analyse statistique.
\end{itemize}

Le matériel de réplication comprend :

\begin{itemize}
    \item les consignes fournies aux participants ;
    \item les tâches expérimentales ;
    \item les tutoriels ;
    \item les questionnaires ;
    \item les environnements d'exécution ;
    \item les tests automatiques ;
    \item les scripts d'analyse ;
    \item les données anonymisées pouvant être rendues publiques.
\end{itemize}


\subsection{Complémentarité avec l'évaluation technique}
\label{subsec:complementarite-evaluations}

L'évaluation technique et l'étude avec participants examinent des
aspects différents de GEAR. Leur complémentarité est présentée dans le
Tableau~\ref{tab:complementarite-publications}.

\begin{table*}[t]
    \centering
    \caption{Complémentarité entre le tool track et l'étude destinée à
    l'article de journal.}
    \label{tab:complementarite-publications}
    \begin{tabularx}{\textwidth}{@{}>{\raggedright\arraybackslash}X>{\raggedright\arraybackslash}X@{}}
        \toprule
        \textbf{Évaluation technique du tool track}
        & \textbf{Étude avec participants du journal} \\
        \midrule
        Couverture des fonctionnalités
            & Capacité des utilisateurs à construire les systèmes \\
        Validité du code généré
            & Temps et effort de développement \\
        Fidélité des traces d'exécution
            & Compréhension et confiance des utilisateurs \\
        Validation des dix connecteurs
            & Modification et maintenance des systèmes \\
        Tests techniques reproductibles
            & Migration entre frameworks \\
        Évaluation automatisée
            & Évaluation contrôlée et entretiens \\
        \bottomrule
    \end{tabularx}
\end{table*}

Ainsi, le tool track démontre que GEAR est techniquement capable de
générer des systèmes corrects pour plusieurs frameworks, tandis que
l'étude avec participants détermine si cette capacité apporte un
bénéfice concret aux développeurs en matière de productivité, de
maintenance, de portabilité et d'expérience utilisateur.